% TOP_PROBLEM: {"id": 30004916, "title": "The Erdős–Rado Sunflower Conjecture", "category_id": 2, "category": {"id": 2, "name": "combinatorics", "display_name": "Combinatorics", "description": "Counting problems, graph theory, discrete structures.", "slug": "combinatorics", "order_index": 2, "created_at": "2026-07-31T15:26:25.671Z"}, "status": "partially_solved"}
% FIRST_POSTED: 2026-09-25
% ENTRY_KIND: statement_only
% !TeX program = lualatex
% Definition and sources only; this file is not an LLM solution attempt.
\documentclass[11pt]{article}
\usepackage[margin=25mm]{geometry}
\usepackage{fontspec}
\setmainfont{Latin Modern Roman}
\usepackage{amsmath,amssymb,mathtools}
\usepackage{unicode-math}
\setmathfont{Latin Modern Math}
\usepackage{xurl,hyperref}
\hypersetup{colorlinks=true,urlcolor=blue}
\setcounter{secnumdepth}{0}
\setlength{\parindent}{0pt}
\setlength{\parskip}{5pt}
\setlength{\emergencystretch}{3em}
\begin{document}
\section*{\texorpdfstring{The Erdős–Rado Sunflower Conjecture}{The Erdős–Rado Sunflower Conjecture}}\label{30004916}
\subsection{Definitions and mathematical statement}
An $r$-petal sunflower is a collection of distinct sets $S_1,\ldots,S_r$ with one fixed core $C$ satisfying $S_i\cap S_j=C$ for every $i\ne j$. Equivalently, their parts outside the core are pairwise disjoint; the core may be empty. The conjecture is
\[
 \forall r\ge3\ \exists C_r>0\ \forall k\ge1:\quad
 |\mathcal F|>C_r^k,\quad|S|\le k\ (S\in\mathcal F)
 \Rightarrow\mathcal F\text{ contains an }r\text{-sunflower}.
\]
The ground set is unrestricted and the exponential base may depend on $r$ but not on $k$.
\par\smallskip\textit{Formulation and concept references:} \href{https://londmathsoc.onlinelibrary.wiley.com/doi/full/10.1112/jlms.70380}{[S1]}.

\subsection{Short English statement}
Must every sufficiently large bounded-size set family contain a sunflower, with a threshold only exponential in the set size?
\subsection{Sources}
\begin{itemize}
\item[S1]Ashwin Rao, The story of sunflowers, Journal of the London Mathematical Society, 2026.
\url{https://londmathsoc.onlinelibrary.wiley.com/doi/full/10.1112/jlms.70380}.
\textit{Formulation source.}
\end{itemize}
\end{document}
